\section{Purpose, scope and status}\label{sec:scope}
This section identifies what the standard defines and how to read its requirements. We consider the Capacitated Vehicle Routing Problem (CVRP) and the Vehicle Routing Problem with Time Windows (VRPTW), with one depot, homogeneous capacity, mandatory single customer visits and static nonnegative arc quantities. The current static families use one matrix as routing cost and, for VRPTW, travel time. We retain each family's objective and units rather than imposing one objective on every benchmark.

The words \must{}, \mustnot{}, \should{} and \may{} express requirements, recommendations and permitted choices within this draft. An implementation claiming a profile \must{} return its specified feasibility verdict and authoritative objective for supported inputs. A resource limit \must{} produce an explicit failure or inconclusive outcome, never a fabricated feasible result. Passing a finite fixture suite is evidence of conformance rather than a proof for every input.

We distinguish three questions. \emph{Operational conformance} asks whether an evaluator implements the prescribed recipe. \emph{Mathematical fidelity} asks whether that recipe computes the intended mathematical model over declared values. \emph{Reproducibility} asks whether recorded artifacts and an environment reproduce the result. A sequential floating-point recipe can satisfy its operational contract while differing from exact summation. An exact rational evaluator can faithfully evaluate rounded matrix entries without reconstructing unrounded geometry. A result reproduced on one host does not establish cross-platform conformance.

The document specifies a future static authority and its alternatives. The preceding investigation of library versions 0.11 and 0.12 remains evidence about those releases, including its earlier compatibility recommendation. It is not silently rewritten as a specification of the new static profiles \autocite{rascoussierNumericalContracts2026}. No Time-Dependent (TD) route-composition or departure-optimization profile is defined here. In particular, this document does not change the historical TD recommendation or claim that integer static costs make TD interpolation exact.

\section{Binary64 values, arithmetic and JSON}\label{sec:binary64}
This section introduces the numerical distinctions needed by the contracts. A binary64 value occupies 64 bits: 1 sign bit, 11 exponent bits and 52 explicit fraction bits. For a normal finite value with sign $s$, exponent field $E$ and fraction field $F$, its exact value is
\[
(-1)^s(1+F2^{-52})2^{E-1023},\qquad 1\leq E\leq2046.
\]
For $E=0$, the finite value is $(-1)^s F2^{-1074}$. This includes subnormal values and signed zero. The all-ones exponent encodes infinities and NaNs. The normal significand has 53 bits of precision \autocite{oracleIEEEArithmetic}.

We write $\RoundBinary(x)$ for rounding an exact real value to the nearest binary64 value, with an exact midpoint rounded to the value whose significand has even least significant bit. We require gradual underflow. In a prescribed binary64 addition, $a\oplus b=\RoundBinary(a+b)$. Rounding after every addition generally differs from rounding only the exact final sum. A compiler or processor that reassociates additions, retains wider unrounded intermediates, flushes subnormals or contracts an expression into a different operation sequence does not implement that prescribed recipe.

Every finite binary64 number is a rational with power-of-2 denominator. Hence we can sum such values exactly and compare exact arrival times without requiring a particular programming language. For example, the binary64 value loaded from \id{0.1} is $3602879701896397/36028797018963968$, whereas the exact decimal value is $1/10$. In Python, \id{Fraction(float("0.1"))} and \id{Fraction("0.1")} construct these different values. Neither operation recovers any physical value that preceded rounding \autocite{pythonFloatingPoint2026,pythonFractions2026}.

The JavaScript Object Notation (JSON) format specifies numeric text, without assigning every consumer a binary64 mathematical interpretation. Its interoperability guidance discusses the commonly exact integer range $[-(2^{53}-1),2^{53}-1]$ \autocite{brayJSON2017}. We therefore declare semantics by field role and profile. Structural integers such as customer indices and route counts remain integers. Numerical quantities such as matrix weights and service times are interpreted according to the selected instance contract, even when a JSON token has no decimal point.

Shortest round-trip output can serialize a binary64 value with far fewer decimal digits than its exact decimal expansion. Under correctly rounded reading, the serialized token reloads to the same finite value. Different spellings can denote that same value, so a textual change can alter a file hash while preserving its interpreted numerical data. Conversely, interpreting the shortest spelling as an exact decimal rational changes its meaning. A loader \mustnot{} infer a decimal contract from the number of printed fractional digits \autocite{pythonFloatingPoint2026}.

The standard loader \must{} retain numeric tokens until their roles and profile are known. It \must{} reject duplicate JSON members, non-JSON numerical extensions such as NaN and Infinity, nonfinite binary64 conversions, and booleans used as numbers. Finite subnormal values are allowed. Values rounded to zero by binary64 conversion have zero numerical meaning under that interpretation. Signed zero is normalized to positive zero for semantic digests and result serialization. Raw artifact bytes retain their original identity. The reference uses token-preserving hooks provided by Python's JSON decoder \autocite{pythonJSON2026}.

\section{Four separate contracts}\label{sec:layers}
This section separates the decisions that together define a benchmark. A \emph{generation contract} records source data, transformations, numerical operations, quantization, random seeds, software and emitted artifacts. An \emph{instance contract} assigns mathematical values and problem structure to those artifacts, including materialized matrices. An \emph{evaluation contract} maps an interpreted instance and submitted routes to a verdict and objective. A \emph{result contract} specifies serialization, comparison, provenance, Best-Known Solution (BKS) storage and the strength of any proof claim.

\begin{center}
\begin{tabularx}{\linewidth}{@{}lXX@{}}
\toprule
Stage & Question & Example obligation\\
\midrule
Generation & How were the published values obtained? & Record the Euclidean recipe and 3-decimal rounding.\\
Instance & What do the values mean after loading or rederivation? & Read each numerical quantity as binary64.\\
Evaluation & What arithmetic and constraints are authoritative? & Use exact sums and exact forward scheduling.\\
Result & What value is ranked and what does a claim establish? & Store a reduced fraction and bind it to the profile.\\
\bottomrule
\end{tabularx}
\end{center}

For this standard, the retained instance interpretation is binary64 after generation or rederivation. This applies to the 3-decimal Poryos/Mamut values as well as full-precision SINTEF-derived matrices. Exact decimal and scaled-grid alternatives are presented explicitly, but they \mustnot{} silently replace the retained binary64 reading. Exact evaluation then means exact arithmetic on those interpreted binary64 values. It does not mean exact decimal arithmetic on their printed spellings.

These choices remain independent of a solver's internal arithmetic. A solver \may{} search with binary64, compensated sums, integers, intervals or rational numbers. Its accepted final incumbent \must{} be checked under the benchmark revision's selected contract. Its lower-bound and optimality claims have the separate obligations in Section~\ref{sec:proof}.
